% !TeX root = ../papers/nc-p3.tex
\section{三割平面の中間座標に対する二次元単純の障害}
\label{sec:trisecant-middle-simple-obstruction}
\begin{theorem}\label{thm:trisecant-middle-simple-obstruction}
前節の二つの異なる有理正規四次曲線を取り、二つの三次因子とその\emph{標準的な基底変更行列}
を連続三座標 $(x,C,y)$ とする。前節の六式とその逆向きの六式を隣接二辺にした
四辺有限の標準正則延長は存在しない。共通二次超曲面が特異でも、交点が非被約でも成立する。
\end{theorem}
\begin{proof}[最後の矢の線型条件]
正則延長を仮定する。隣接二辺は六独立、三次交叉は次元4であり、
前節の証明により第四文字の点座標空間を
\[
 \mathcal G=\{G\in\mathcal H:T^{\mathsf t}GT\in\mathcal H\}
\]
に同定できる。$G(k,h)=\phi_G(kh)$、
$G(Tk,Th)=\psi_G(kh)$ として $\phi_G,\psi_G\in S_8^*$ を取る。
$u_0=s,u_1=t$ とすると四テンソルは中間位置の $ij$ 成分で
$G(T(u_i x),u_j y)$ である。従って $(y,G)$ 辺の関係は
$\phi_G(fy)=0$ ($f\in S_5$)、$(G,x)$ 辺は $\psi_G(fx)=0$ ($f\in S_5$) で与えられる。
収縮で $S_4 S_1=S_5$ が張られるため全関係が得られる。

$e_0,e_1,e_2=s^2,st,t^2$ とし、未知数 $G_{ij}\in\mathcal G$、$0\le i,j\le2$ に
\[
 \sum_j\phi_{G_{ij}}(e_j f)=0\quad(0\le i\le2,\ f\in S_6),\qquad
 \sum_i\psi_{G_{ij}}(e_i f)=0\quad(0\le j\le2,\ f\in S_6)
\]
を課す（$i,j$ はそれぞれ三添字、$f$ は六次式）。三十六未知数への四十二条件であるが、
解は必ず非零である。実際 $a,b\in S_2$ に対し、第一群に $f_i=T(e_i a)b$、
第二群に $f_j=a T^{-1}(e_j b)$ を使って総和を取ると両辺は同じになる。
共に $\sum_{ij}G_{ij}(T(e_i a),e_j b)$ だからである。少なくとも七本が独立な従属性となる。
第一群だけに出る係数への写像
$a\otimes b\mapsto(T(e_i a)b)_i$ の核を考える。
核元を第二変数の点 $v\in\mathbb P^1$ で評価して得る二次式列 $a(v)$ は、
二項四次汎関数 $T^*\operatorname{ev}_v$ の中央三行触媒行列
$(T^*\operatorname{ev}_v(e_i e_j))_{ij}$ の核に入る。
一般階数が1なら点曲線 $T^*R$ は $R$ と等しくなる。実際階数1の四次Hankel汎関数は
一点評価であり、これは二次Veronese曲線の方程式または単項式の連比で確認できる。
よって一般階数は2以上。階数3なら写像の核は零、2なら核列は原始な同次列の多項式倍である。
その列の次数は1以上である。定数二次式 $a$ なら $T(e_i a)$ が全 $v$ で零となるからである。
二次の列に入る倍数の空間は従って二次元以下で、$9-2=7$ の下界を得る。
\end{proof}

\begin{proof}[単純の構成と両向きの頂点拡張]
非零解を取り、四頂点の空間を矢位置 $x,C,y,G$ の順に
$S_2,k^2,k^2,S_2^*$ として一周させる。
二項三次式にまとめた $x$ の矢は $p\in S_2$ を
$(t p,-s p)^{\mathsf t}$ に送る（各三次係数が一つの矢）。
$C_{ij}$ の矢は二番目の $k^2$ から見た行番号 $j$、初めの $k^2$ から見た列番号 $i$ の行列単位とする。
$y$ の矢は、同じ $p\mapsto(t p,-s p)^{\mathsf t}$ の三次係数の転置である。
$G$ 位置の矢 $l\in\mathcal G^*$ は行列 $(l(G_{ij}))_{ij}:S_2^*\to S_2$ である。
初めの辺の関係は $s(tp)+t(-sp)=0$ に $T$ を作用させて $s$ または $t$ を掛けたものになる。
次の辺は転置と $T^{-1}$ で同じである。残りの二辺は先の線型条件である。
$S_5$ に $s,t$ を掛けたものが $S_6$ を張るからである。

頂点0を最後の全行列の像の包に制限し、頂点3を共通核で割る。その次元を $r,b$ とする。
両方正である。得る表現は単純である。非零な $p$ に対し $tp,-sp$ は独立なので
$x$ の四矢の像は次の二空間を張る。中間の四行列単位も同じ性質を持つ。
$y$ の四矢は非零 $k^2$ ベクトル上で元の $S_2^*$ 全体を張る。
双対の核が非零なら二次式と非零一次式の積が零になるためである。
共通核を除いた最後の位置では非零ベクトルを全矢が消すこともない。
従って非零部分表示は二つの中間空間を全て含み、前後の空間も全て含む。

定理\ref{thm:generic-corner-small-simples}から $r=b=2$ でなければならない。
しかしこの場合には元の三次元 $S_2$ の残る一方向を頂点0に戻すと頂点単純を商とする
非分裂拡張になる。戻した方向の $x$ の作用は元の二方向のどの線型結合とも一致しない。
逆に頂点3で割った一方向を戻すと頂点単純を部分とする非分裂拡張になる。
$y$ の元の全像が三空間を張るため補空間で分裂できない。
いずれも他の矢は包含または商との合成にすれば全関係を保つ。
これは同じ定理の二次元均衡単純での片向き消滅に反する。
\end{proof}

\begin{corollary}\label{cor:generic-seven-cofactor-pair-excluded}
七余因子の行列正規形における一般のパラメータ $T$ では、標準中間座標以外を許しても
この標識付き反対配置の対を含む有限な標準正則窓はない。
\end{corollary}
\begin{proof}
一般の $T$ では標準中間辺の両射影は有限で、いずれも階数3の点のグラフとなる。
両射影の有限性は付属計算の有理標本で、一つの平面上の十五個の四次最大小行列式が
全四次式を張ることから開条件で確認できる。階数2以下があれば反対射影に直線が出る。
また両側の曲線を入れ替えた逆の標準中間辺でも同じ有限性を検査できる。
長さ20の共通三割平面から中間配置への射が同型になる。
この中間配置では八次式の全二階偏導の六次商への評価が単射であることを、
同じ標本において良い素数で確かめる。従って三重消滅八次式はない。

標本では中間配置が被約で、いずれの四点も独立であることを次のように確かめる。
四次小行列式で割った商の三次と四次は二十次元であり、$C_{00}$ 倍は同型である。
これで正規化した $C_{01}$ 倍の有理数上の特性多項式は次数20で、
その原始整数式を $31,53,61$ で還元したとき、全て平方自由で
既約因子次数はそれぞれ $(20),(1,19),(1,2,17)$ である。
従って配置は被約でこの比が全点を分離する。Galois群は推移的であり、19巡回を含むので
二重推移的である。最後の巡回型の17乗は互換であるから、共役で全ての互換を含む。
このFrobenius巡回型判定は先頭係数と判別式が単元である各素数で、根に対する
剰余体のFrobenius置換を分解群に持ち上げることによる。
全点が平面を張るだけなら三次評価が単射でないので、独立な四点があり、
対称群の作用で全ての四点が独立である。有限エタール族で相異なる四点の順序組も有限エタールだから、
これらの条件は一般の $T$ に保たれる。

実際の有限窓で $x$ の隣にある座標を $V$ とする。
$(x,C),(x,V)$ のいずれも同じ標識付き長さ20配置への有限完全交叉グラフである。
有限辺の跡による直交性（定理\ref{thm:finite-square-chart}）から、
$C^2,V^2$ はともに $x^2$ の直交空間の単元倍である。従って
$C$ の十個の二次式を基底に持つ対称行列 $M$ があり、各点で $V$ の平方に比例する。
定理\ref{thm:opposite-root-triple-octics}の被約の場合の行列計算は、
この行列の性質、$C$ の三次評価の単射性、$C$ の直線上に四点がないことだけを使う。
従って $V$ と $C$ が標識付き射影同値でなければ $C$ を三重に通る八次式があり矛盾する。
$(x,V,y)$ の両辺は標準の二辺と同じ配置グラフを持ち、各グラフを通る双一次式は
六個の定義式で尽きる（完全交叉のKoszul分解を $(1,1)$ 捻る）。
定理\ref{thm:trisecant-middle-simple-obstruction}を適用する。
\end{proof}

\path{computations/verify_trisecant_middle_simple_obstruction.py} はこの三素数での
因子次数と八次評価を再現する。また、共通の階数5二次超曲面を保つ別の有理標本で、
四辺の有限性と有限Frobenius条件、反対階数16を検査する。
四十二条件の核は有理数上で四次元で、得られる解の一つでは $(r,b)=(3,3)$ である。
この標本では次元ベクトル $(3,2,2,3)$ だけで矛盾が確認できる。
